% Functional diagram: orthogonal (straight-line) layout on a 4-column grid.
% Columns: inputs (x=0), functions F1/F2 + state (x=3.4), functions F3/F4 (x=6.8), sinks (x=10.4).
\begin{figure}[htbp]
\centering
\scalebox{0.9}{%
\begin{tikzpicture}[x=1cm,y=1cm,
  lbl/.style={font=\scriptsize,fill=white,inner sep=1.5pt},
  dot/.style={circle,fill=black!65,inner sep=1.3pt}]

% ---------------------------------------------------------------- nodes
\node[funcbox,fill=yellow!25,text width=1.9cm,minimum height=9mm] (sun) at (0,2.6)
  {Solar\\ irradiance};
\node[funcbox,text width=2.4cm,minimum height=11mm] (f12) at (3.4,3.7)
  {\textbf{F1} Reflect short-wave radiation\\[1pt] \textbf{F2} Shade a surface};
\node[funcbox,fill=red!8,text width=2.4cm,minimum height=11mm] (ts) at (3.4,1.4)
  {Absorbed heat\\[1pt] \itshape state: surface temperature $T_s$};
\node[funcbox,text width=2.4cm,minimum height=9mm] (f3) at (6.8,3.7)
  {\textbf{F3} Dissipate / transfer heat};
\node[funcbox,text width=2.4cm,minimum height=9mm] (f4) at (6.8,0.5)
  {\textbf{F4} Evaporate liquid};
\node[funcbox,fill=zhawblue!10,text width=1.9cm,minimum height=9mm] (wat) at (6.8,-1.1)
  {Water in};
\node[funcbox,fill=gray!12,text width=2.0cm,minimum height=9mm] (sky) at (10.4,3.7) {Sky};
\node[funcbox,fill=gray!12,text width=2.0cm,minimum height=9mm] (air) at (10.4,2.3) {Street-canyon air};
\node[funcbox,fill=gray!12,text width=2.0cm,minimum height=9mm] (vap) at (10.4,0.5) {Vapour to air};
\node[funcbox,fill=gray!12,text width=2.0cm,minimum height=9mm] (bld) at (10.4,-2.1) {Building interior};

% ---------------------------------------------------------- energy flows
% sun -> F1/F2 and -> absorbed heat (branch b1)
\coordinate (b1) at (1.7,2.6);
\draw[flowarr] (sun.east) -- (b1) |- (f12.west);
\draw[flowarr] (b1) |- (ts.west);
\node[dot] at (b1) {};
\node[lbl,anchor=east] at (1.58,2.0) {absorbed};

% reflected: over the top, straight down into the sky
\draw[flowarr] (f12.north) -- ++(0,1.1) -| (sky.north);
\node[lbl] at (6.9,5.35) {reflected};

% absorbed heat -> F3 and F4 (branch b2)
\coordinate (b2) at (5.0,1.4);
\draw[flowarr] (ts.east) -- (b2) |- (f3.west);
\draw[flowarr] (b2) |- (f4.west);
\node[dot] at (b2) {};

% F3 -> sky (radiation) and -> canyon air (convection), branch b3
\coordinate (b3) at (8.6,3.7);
\draw[flowarr] (f3.east) -- (b3) -- (sky.west);
\draw[flowarr] (b3) |- (air.west);
\node[dot] at (b3) {};
\node[lbl] at (8.7,4.0) {radiation};
\node[lbl,rotate=90] at (8.6,3.0) {convection};

% residual: absorbed heat -> building interior (to be minimised)
\draw[flowarr,draw=red!70] (ts.south) |- (bld.west);
\node[lbl,text=red!70] at (6.8,-1.82) {residual heat into the building};

% -------------------------------------------------------- material flows
\draw[wateharr] (wat.north) -- (f4.south);
\draw[wateharr] (f4.east) -- (vap.west);

% ------------------------------------------------------ information flow
% T_s drives F3 and F4: orange dashed lines run parallel to the energy paths
\draw[infoarr,transform canvas={xshift=3.5pt,yshift=-3.5pt}] (ts.east) -- (b2) |- (f3.west);
\draw[infoarr,transform canvas={xshift=3.5pt,yshift=-3.5pt}] (b2) |- (f4.west);
\node[lbl,text=orange!85!black,font=\scriptsize\bfseries] at (5.85,2.55) {$T_s$ signal};

% ------------------------------------------------------- system boundary
\begin{scope}[on background layer]
  \node[draw=zhawblue,dashed,rounded corners=4pt,fit=(f12)(ts)(f3)(f4),
        inner sep=9pt] (fb) {};
\end{scope}
\node[font=\scriptsize\bfseries,color=zhawblue,anchor=east]
  at ($(fb.north east)+(-0.1,0.28)$) {SYSTEM BOUNDARY};

% ---------------------------------------------------------------- legend
\node[font=\scriptsize,anchor=west,align=left] at (-1.1,-3.1)
  {\textcolor{black!65}{\textbf{---}} energy flow \quad
   \textcolor{zhawblue}{\textbf{--\,--}} material flow (water, vapour) \quad
   \textcolor{orange!85!black}{\textbf{--\,--}} information flow \quad
   \textcolor{red!70}{\textbf{---}} flow to be minimised};
\end{tikzpicture}}
\caption{Solution-neutral functional diagram; the system boundary is the outer layer of the building. Solid arrows are energy flows, blue
dashed arrows the material flow of F4 (water in, vapour out), the orange dashed
arrows the information flow: the surface temperature drives the strength of F3 and
F4. F5 (resist degradation) acts on everything inside the boundary and is not
drawn.}
\label{fig:funcdiagram}
\end{figure}
